\documentclass{article}
\title{ノートと読み書きの技法 — 頭のまま残り続ける、4つの道具}
\author{TeX 図書館}
\lstdefinelanguage{Text}{basicstyle=\ttfamily,breaklines=true,columns=fullflexible}

\begin{document}

\section{はじめに — ノートは記憶の足場}

\textbf{ノートを取ったのに頭に残らない}、という悩みから、人はノートをとっても、残りません。本書は、\textbf{4 つの道具}で、この悩みを解きます。

\subsection{ノート技法が効く理由}

記憶は 1 つの箱ではありません。\textbf{符号化} \(\rightarrow\) \textbf{保存} \(\rightarrow\) \textbf{検索} という 3 段階の工事として考えると、ノートという行為が何をしているのかがはっきりします。

\begin{itemize}
\item \textbf{符号化}: 見たものを、教えられる形に変える(日本語にする、絵にする、質問に変える)
\item \textbf{保存}: 変換されたものを、時間とともに保持する
\item \textbf{検索}: 必要なときに\textbf{取り出す経路}があること
\end{itemize}

ノートという行為は、このうち \textbf{符号化} と \textbf{検索の経路をつくる} 仕事です。ノートの価値は、\textbf{あとで自分で辿れる道を作ること}にあります。

\subsection{効果量という目で見る}

科学のほうは、この手の話を「効いた / 効かない」で語らず、\textbf{効果量} \(d\) で語ります。
\[ d = \frac{\bar{x}_1 - \bar{x}_2}{s} \]
\(d = 0.2\) が小さい効果、\(d = 0.5\) が中くらい、\(d = 0.8\) が大きい効果、というのが目安です。\textbf{『統計学入門』第 10 節}で読むこの数字を、ノート術の議論にもそのまま使えます。

ノートに関する研究の多くは、効果量が小さいことを認めています。量ではなく、設計の質に左右されるという結果のほうが、技法の設計の質を示しています。

\begin{quote}
\textbf{【補足】}\ 効果量の目安: \(d = 0.2\) 弱い、\(d = 0.5\) 中くらい、\(d = 0.8\) 強い。\textbf{『統計学入門』第 10 節}の「効果量」を、本書でも使い回してください。
\end{quote}

\subsection{記憶の代わりではなく、記憶の足場}

\begin{quote}
\textbf{【定義】}\ ノートとは、\textbf{記憶そのものではなく、記憶を手繰る足場}のこと。あとから自分で引き出せるようにするための、\textbf{外部に置いてから引く道具}にすぎない。\end{quote}

この言い方には 2 つの実際的な意味があります。\textbf{ノートを残したこと}を、\textbf{覚えたこと}の代わりにしない。そして、\textbf{『総合問題集』}がやっている「解けなかった問題は 1 週間後に再解く(間隔反復)」は、まさにこの足場を使う定式です。\textbf{ノートを作って終わり}ではなく、\textbf{ノートを引いて続けて初めて記憶になる}。

\begin{quote}
\textbf{【演習】}\ あなたは 1 冊の本を読んでノートを作った。それは「記憶の足場」になっているか。判定の観点を 3 つ挙げよ。
\end{quote}

\textbf{解答の指針:} 3 つの観点は「\textbf{あとで自分の言葉で説明できるか}」「\textbf{索引を引いて 1 分で辿れるか}」「\textbf{日付を隔てて出し直して答えられるか}」。どれか 1 つでも欠ければ、足場ではなく写しの捨て場です。

\section{4 つの道具の全容}

この本ではノート技法を 4 つの道具に分けます。道具ごとに性格が違うので、\textbf{全部使う}のではなく、\textbf{何のためにどれを出すか}で使い分けます。

\subsection{道具 1・2・3・4 は何が違うのか}

\begin{center}
\begin{tabular}{|l|l|l|}
\hline
道具 & やること & 主な働き \\
\hline
1 書き写す & 頭の中のものを外部に置く & 符号化 \\
\hline
2 結びつける & 既存の記憶に結びつける & 検索 \\
\hline
3 作り変える & 自分の形に変える & 符号化 \\
\hline
4 引き出す & 隠して引き出そうとする & 検索 \\
\hline
\end{tabular}
\end{center}

道具 1 と 3 はどちらも符号化方向、道具 2 と 4 はどちらも検索方向です。違いは、1 が「\textbf{浅く写す}」こと、3 が「\textbf{深く変形}」すること、2 が「\textbf{道を張る}」こと、4 が「\textbf{張った道を実際に歩く}」ことです。

\subsection{使い分けの基準 — どの場面で使うか}

同じ「ノート」でも、目的が変われば選ぶ道具が変わります。

\begin{center}
\begin{tabular}{|l|l|}
\hline
場面 & 主役と理由 \\
\hline
政治の年表 & 書き写す(1)。事実を並べるだけで成立する \\
\hline
化学反応式 & 作り変える(3)。電子の流れを説明できなければ組めない \\
\hline
物理的名詞 & 結びつける(2)。どの定理から来たかが要る \\
\hline
試験準備 & 引き出す(4)。見返しても点が上がらない \\
\hline
\end{tabular}
\end{center}

政治の年表は事実を並べればよく、切り詰めたら壊れます。化学反応式は「なぜそう書くか」を自分の言葉で説明できなければ、試験で再現できません。物理の精密な用語は、どの理論のどの定理から来たかに意味があります。\textbf{道具の順番は固定していません}。場面の順序に従います。

\subsection{学習の軸 — 収集・構造・変換・出力}

本書の 4 道具は、別の分類法ではなく、\textbf{情報処理の 4 段階}に整理してあります。

\begin{center}
\begin{tabular}{|l|l|l|}
\hline
軸 & 道具 & 得られるもの \\
\hline
収集 & 1 書き写す & 材料 \\
\hline
構造 & 2 結びつける & 材料どうしの関係 \\
\hline
変換 & 3 作り変える & 自分の中の理解 \\
\hline
出力 & 4 引き出す & 使える能力 \\
\hline
\end{tabular}
\end{center}

「\textbf{触ったのに残らない}」という症状の多くは、\textbf{収集の段階で止まっている}ことにあります。ノートが「集まっただけ」なら、それはまだ収集です。この 4 軸のままでも、どの段階が止まっているかは診断できます。

\section{道具の前提 — 記憶の仕組み}

道具を使い分ける前に、\textbf{道具が作用する仕組み}を 4 つ押さえます。すべて既存の本から持ってきた材料です。

\subsection{記憶は 3 段階の工事である}

先に述べたとおり、記憶は符号化・保存・検索の 3 段階です。ここで重要なのは、\textbf{検索の経路がないと、保存したものは取り出せない}という一点です。

\textbf{『大学数学入門』}で「固有値を計算する方法」を覚えたとします。しかし「\textbf{固有値を使う場面}」と結んだ瞬間、同じ概念が別の文脈でも検索できるようになりました。\textbf{検索の道がない知識は、保存しても取り出すときに使えません}。これが\textbf{孤立}の問題です。

\begin{quote}
\textbf{【定義】}\ \textbf{符号化深度}(depth of processing)とは、情報が多大まで処理されたかのこと。\textbf{浅い符号化}は音や形の観察(見えた文字をそのまま写す)で、\textbf{深い符号化}は意味づけ・例化・関係づけ(自分の言葉に直す、類比をかける)です。
\end{quote}

\subsection{作業記憶の容器の広さ}

Miller(1956) は、\textbf{作業記憶の容量}が「\(7 \pm 2\)」個だと報告しました。2010 年になって Cunningham と Cowan が再検討し、「\(4 \pm 1\)」個ではないかと議論されています。

\begin{quote}
\textbf{【補足】}\ 「\(7 \pm 2\)」と「\(4 \pm 1\)」のどちらを採っても結論は同じで、\textbf{作業記憶に同時に載せられるのは数個}までです。よってノートは「\textbf{1 画面=} 15 行以内」が現実的です。
\end{quote}

ここで数学で裏を取ります。状態数 \(n\) の行列の累乗を素朴に計算すると、乗算回数は \(O(n^3)\) です。16 状態なら \(16^3 = 4096\) 回。\textbf{対角化} \(A = PDP^{-1}\) を使えば、\(A^{16} = PD^{16}P^{-1}\) となり、対角部分の累乗は成分ごとに計算できます。\textbf{ノートで圧縮すること}は、\textbf{作業記憶の容量制限を回避する技術}でもあるのです。

\subsection{注意と認知負荷 — 3 つの load}

ノートを取る瞬間に作用するのが\textbf{認知負荷}(cognitive load)です。Sweller の分類では 3 種類に分けられます。

\begin{center}
\begin{tabular}{|l|l|l|}
\hline
種類 & 内容 & 減らし方 \\
\hline
内在的負荷 & 材料そのものの難しさ & 材料を分解する \\
\hline
外在的負荷 & 書き方の体裁や量 & 書く量を減らす \\
\hline
付随する負荷 & 例や図など理解を支えるもの & 适量増やす \\
\hline
\end{tabular}
\end{center}

内在的負荷は単純に減らせません。それが内容の難易度だからです。減らせるのは外在的負荷——「ノートをきれいにすること」「書き写すこと」——のほうです。

\begin{quote}
\textbf{【補足】}\ 注意とは、3 つの負荷の中で唯一、\textbf{自分で増やせて自分で減らせるもの}。ノート術は、\textbf{頭を使う場所を増やす}技術になります。
\end{quote}

\subsection{ノートの前に問うべきこと}

本書を貫く最初の行動ルールです。\textbf{ノートを取る}前に、\textbf{この情報を後でどう使うか}を 1 文で書く。

\begin{center}
\begin{tabular}{|l|l|}
\hline
後から使う場面 & そのとき必要なノートの形 \\
\hline
定義を説明するとき & 定義と 1 例と 1 つの非例 \\
\hline
手順を扱うとき & 手順と注意点 \\
\hline
試験で答えるとき & 自分の理屈と、証拠となる数値 \\
\hline
\end{tabular}
\end{center}

この表のように、\textbf{使う側が決まれば書くべき形も決まる}。逆に言えば、後での場面が決まっていない情報は、ノートに蓄積しても出てきません。\textbf{Cornell ノート方式}は「右半分に自分のメモを書く」という方向の試みですが、列挙された形がそのまま「記憶の足場の型」になるとは限りません。\textbf{使う場面から、型を書く}。これが本書の順序です。

\section{道具 1: 書き写す — 材料を外部に置く}

\subsection{消えるノートと残るノート}

最も密度の低いノートは「\textbf{全部そのまま写した}」ノートです。最も密度の高いノートは「\textbf{自分で言い換えた}」ノートです。見た目はどちらも「文字が並んでいるだけ」ですが、中身の密度が違います。

\begin{center}
\begin{tabular}{|l|l|l|}
\hline
ノート & 1 週間後 & 1 年後 \\
\hline
暗記したノート & 半分消える & 痕跡がない \\
\hline
自分で言い換えたノート & 大半が残る & 拡張に使える \\
\hline
\end{tabular}
\end{center}

下地にあるのは、\textbf{書き写すときの処理の深さ}です。\textbf{文字をそのまま写す}のは、見えたものを機械的に移すだけの浅い符号化です。\textbf{写しながら「これはつまり何だ」と考える}のは、深い符号化です。\textbf{耳で聞いただけで消えるノート}とは、\textbf{意味づけを通っていないノート}のことです。

\subsection{手書きとタイピングの論争}

「\textbf{手で書くか、キーボードで打つか}」の議論を、実測値つきでまとめます。

\textbf{Mueller \& Oppenheimer(2014)} は、ノートをすべてタイピングした学習者が、\textbf{文章脈のある質問}(指示された考えを組み立て、関係を述べる問題)で手書き群を下回ったと報告しました。理由は、\textbf{タイピングが速すぎて写せてしまう}ことです。単語数を数えると、手書き群は平均 156 語、タイピング群は平均 261 語で、差は約 1.1 標準偏差(\(d \approx 1.1\))に及びます。\textbf{同じ内容の余計な量}に見えて、その逐語の多さそのものが効果を下げた、というのが解釈です。

\begin{quote}
\textbf{【注意】}\ ただし「\textbf{タイピングが悪い}」と読むのは早計です。Bui, Myerson \& Hale(2013)は、\textbf{要約しながら打つ}条件ではタイピングが手書きを上回りました。\textbf{道具ではなく、書き方が効いている}のです。
\end{quote}

\begin{quote}
\textbf{【補足】}\ 手書きとタイピングの差は、\textbf{書いているときの処理の深さ}が主因。\textbf{丸写しでも、手書きでも同じように薄くなります}。
\end{quote}

\subsection{よく落とす書き写し}

\begin{enumerate}
\item \textbf{丸写し}: 黒板の字形をそのまま写す。あとで「\textbf{誰のどの発言か}」が分からない
\item \textbf{2 回読む}: 1 回目は「理解したつもり」で読む。2 回目に抽象化は起きない
\item \textbf{全部に取る}: 全部書き写す。\textbf{重要度}でサイズを変える(重要なら大きく、普通なら 1 行、蛇足はなし)
\item \textbf{手書きのノートをデジタルにしない}: 手で書いた方が検索経路が早いのは、\textbf{位置}そのものが目印になるからです
\end{enumerate}

\begin{quote}
\textbf{【定義】}\ \textbf{外部保存効果}(external storage effect)とは、\textbf{ノートに書かなかったものは記憶しなくてよい}という原則。ノートを取る人の労力と省略を同時に防ぐ、生命の根拠です。
\end{quote}

\begin{quote}
\textbf{【演習】}\ 90 分の講義でノートを取る。\textbf{書き写しの量}を最適化するため、作業記憶の容量 \(4 \pm 1\) を使って要約の長さを決めよ。何行に収めるべきか、理由とあわせて述べよ。
\end{quote}

\textbf{解答の指針:} 常に \(4 \pm 1\) 個のチャンクしか作業記憶に同時に載せられません。よって\textbf{要約は 4 行以内}に収めます。\textbf{4 行(結論 1 行、用語 1 行、数値 1 行、疑問 1 行)}だけ書き、残りは元の情報源に委ねる、という運用が実務的です。\textbf{検索練習で再現するとき}、この 4 行だけで答えが出なければ、要約の粒度が粗すぎるままです。

\section{道具 1 の実装 — 自分の雛形を持つ}

\subsection{3 つの雛形}

道具 1 の実装は、\textbf{使う場面ごとに雛形を 3 つだけ持つ}ことです。

\begin{itemize}
\item \textbf{Daily note(記録ノート)}: 1 日 1 ファイル。会議・接触・考えの芽。\textbf{時系列で浅く}書いてよい
\item \textbf{文献ノート}: 1 冊 1 ファイル。\textbf{その本から得た主張}だけを、自分の言葉で
\item \textbf{参考文献ノート}: 「\textbf{どの主張を、どの本から持ってきたか}」の台帳
\end{itemize}

3 つに分けるのは、\textbf{自分の考え}と\textbf{他人の主張}を混ぜないためと、PARA 系(Projects / Areas / Resources / Archive)の Resources 箱に対応させるためです。\textbf{ファイル数を増やさない}。

\subsection{Daily note の雛形}

\begin{lstlisting}[language=Text]
# 2026-09-29

## 今日の結論(1 行で)
- 確率分布の密度関数は「面積が 1 の山」であり、縦軸は確率密度であって確率ではない。

## 残したノート
- 『統計学入門』§6: 正規分布は平均と標準偏差の 2 つで決まる。標準化変数 z は
  「平均から標準偏差何個分か」を表す量。

## 引いてみる(書かずに、小声で)
- z = 2 の片側確率は約 2.3%。この数字はなぜそうなるか。-> 描ける。

## 不明点(後で調べる)
- 標本標準偏差の分母に n ではなく n-1 を使う理由は?(不偏性)
\end{lstlisting}

この雛形で重要なのは、\textbf{「引いてみる」と「不明点」の 2 ブロック}です。\textbf{道具 1 は道具 3 と道具 4 の入口を同時に開いている}。1 ファイルの中に、既に検索練習と「\textbf{自分の言葉}」への道が埋め込まれています。

\subsection{参考文献の書き方}

参考文献ノートは、\textbf{道具 1 と道具 2 の組み合わせ}として費用対効果が高いものです。書式は 1 行 1 件。\textbf{主張}が左、\textbf{出所}が右。

\begin{lstlisting}[language=Text]
主張: p 値は「H0 のもとで観測結果以上に極端になる確率」で、
      H0 の正しさの確率ではない。
出所: 『統計学入門』§9
読了: 2026-09-20
自分の言葉: 標本を取り直せたかどうかではなく、この H0 の下で
            もういちど出ている見込みの確率。

主張: Miller(1956) の 7±2 は当時の評価であり、
      Cowan(2001) が 4±1 への再検討を行った。
出所: 記憶容量の古典論文(Miller 1956 / Cowan 2001)
読了: 2026-09-25
自分の言葉: 正確な値は議論の対象だが、「数個」という感覚だけ
            押さえておけば十分。
\end{lstlisting}

\textbf{「自分の言葉」を必ず 1 文書く}。これが\textbf{道具 3 の実装}への足場であり、丸写しの再発を 1 文で防ぎます。

\begin{quote}
\textbf{【演習】}\ 昨日読んだ 1 冊の\textbf{参考文献ノート}を、上の雛形に合わせて\textbf{2 分で 3 件だけ}書いてみよ。書けない場合、それは\textbf{読んでいない}ので、どこに戻すべきか述べよ。
\end{quote}

\textbf{解答の指針:} 「自分の言葉」が書けない 1 件は、\textbf{その 1 節だけ} 2 行(結論 1 行 + 不明点 1 行)に退避させます。\textbf{退避した 1 行}が、そのまま次の白紙法の日付のノートになります。これが「\textbf{集まっただけ}」を\textbf{消化済み}に変える最小のアクションです。

\section{道具 2: 結びつける — 記憶を網にする}

\subsection{孤立した知識は消える}

孤立した知識は消える。理由は符号化ではない。\textbf{検索の経路が 1 本もない}からです。

Ebbinghaus の忘却曲線は、時間とともに記憶が衰えることを示します。保持率の減衰は、おおよそ指数関数の形
\[ R(t) \propto e^{-t/S} \]
で表せます。ここで \(t\) は経過時間、\(S\) は\textbf{安定性}の指標です。曲線の傾きが緩やかになるのは、\textbf{記憶が強化されている}からにほかなりません。

重要なのは、\textbf{符号化が深いほど、あるいは他の概念と結ばれているほど \(S\) が大きくなる}という点です。\textbf{孤立したノートは \(S\) が小さく、消えやすい}。\textbf{他と結びついたノートは \(S\) が大きく、残りやすい}。

\subsection{概念マップの書き方}

\textbf{概念マップ}とは、\textbf{概念を頂点、関係を辺}として描いたグラフです。\textbf{ノート同士を「A である」以外で結ぶ}のが本体です。

\begin{enumerate}
\item \textbf{中央にテーマ}を書く(1 つのノートなら、その中心語)
\item \textbf{主枝を 3〜5 本}書く(根拠は 3〜5 本。それ以上は作業記憶の容量を超える)
\item \textbf{各枝に概念}を書く(1 枝 1 概念。1 行が理想)
\item \textbf{枝同士を線で結ぶ}(ここが map であり mind map との差)
\end{enumerate}

\textbf{4 の「枝同士を結ぶ」が Mind map と Concept map の境界}です。\textbf{Mind map} は 1 本のテーマから放射状に伸ばすだけで、枝同士の関係を書きません。\textbf{Concept map} は「\textbf{分散と統合}」という関係も書き加えます。

\begin{lstlisting}[language=Text]
                        [ 中央: p 値 ]
            /            |              \
   (内在的説明)      (外在的説明)      (誤用)
         |                |                |
   「H0 のもとで     「H0 は正しい      「H0 の正しさの
    観測以上に極端      確率」→ 誤り     確率を p で表す」
    になる確率」         (方向が逆)       (意味が違う)
         |                |
     (検出力と       (事前確率と
      関係)          関係)
\end{lstlisting}

\subsection{バックリンク — 片側リンクに価値がある}

\textbf{Obsidian} など多くのノートツールでは、あるノートから別のノートへ張ったリンクは、自動で「\textbf{戻りリンク(バックリンク)}」として返されます。これが Zettelkasten の真価です。\textbf{カード 1 枚から、全体の姿が見える}。

\begin{center}
\begin{tabular}{|l|l|}
\hline
リンクの種類 & 見えるもの \\
\hline
片側リンク(前方向) & 自分が向かっている先 \\
\hline
バックリンク(戻り) & 「誰が自分を引いたか」が自動で見える \\
\hline
\end{tabular}
\end{center}

\begin{quote}
\textbf{【定義】}\ \textbf{Zettelkasten}(ツェッテルカステン)は、社会学者 Niklas Luhmann のカード法です。\textbf{1 枚のカードは 1 つの主張だけ}を持ち、\textbf{カード ID を振り}、\textbf{新しいカードから既存のカードへリンク}します。全体がフラットな番号の集合になり、後から「ある概念から始まる物語」を任意の切り口でたどれます。
\end{quote}

\begin{quote}
\textbf{【定義】}\ \textbf{原子ノート}(atomic note)とは、\textbf{1 枚のノートが 1 つの主張だけ}を持つこと。2 つの話題を 1 枚に混ぜないことが要です。
\end{quote}

\subsection{1 例の概念マップ(ASCII)}

以下は、『統計学入門》第 9 節(仮説検定)を実際に 1 周したときのマップです。

\begin{lstlisting}[language=Text]
              [ 仮説検定 ]
        /         |          \
   [ p 値 ]  [ 有意水準 ]   [ 検出力 ]
      |           |            |
   ┌──┴──┐    答えの限界      効果量
   |     |                      |
 内在  外在                 標本数 n
 説明  説明(誤り)                |
   |     |                  ┌─────┘
 p 値の  p 値の誤用          └→ 効果量 d
 定義
\end{lstlisting}

\section{道具 2 の実装 — カードを設計する}

\subsection{Zettel カードの設計}

原子ノートの実装は、\textbf{4 つの枠}でカードが収まるようにすることです。\textbf{ID・日付・本文・リンク}の 4 枠です。

\begin{lstlisting}[language=Text]
ID:       20260929-0310
見出し:   仮説検定の p 値は「H0 の下の確率」ではない
本文:     p 値は帰無仮説 H0 が真のもとで、観測結果以上に極端な
          データが得られる確率。H0 の正しさの確率ではない。
          方向が逆である。検出力と効果量と組み合わさって
          はじめて意味を持つ。
リンク:   [[p 値]] [[仮説検定]] [[検出力]] [[効果量]]
質問:     なぜ p = 0.04 なのに「有意でない」と言った人がいるか?
日付:     2026-09-29
\end{lstlisting}

\begin{quote}
\textbf{【補足】}\ \textbf{質問}欄は、\textbf{理解度の目印} であり、\textbf{検索練習のカード} でもあります。\textbf{自分から出した質問}を 1 つ書くだけで、引き出す回路が機能しはじめます。
\end{quote}

\subsection{バックリンクで「引き出す 10 分」}

\texttt{[[ノート名]]} というリンクを 2 つ張るだけで、バックリンクは自動的に返ってきます。

\begin{enumerate}
\item \textbf{1 分目}: \textbf{検索}で 1 日分のファイルを 1 つだけ選ぶ(新旧は問いません)
\item \textbf{2〜3 分目}: そのファイル\textbf{だけ}を開く。見返さないので、\textbf{設問}を 1 つ選んで読む
\item \textbf{4〜6 分目}: \textbf{白紙}(A4) を 1 枚取り、\textbf{カードの本文}を読み\textbf{自分の言葉}で書き写す。あわせて、\textbf{リンク先}を 1 つだけ書く
\item \textbf{7〜8 分目}: 逆にたどって、\textbf{リンク先からこのカードに戻れるか}確かめる
\item \textbf{9〜10 分目}: \textbf{?} に答えられれば終了、できなければ \textbf{「不明点」に追記}して終了
\end{enumerate}

\begin{quote}
\textbf{【演習】}\ \textbf{白紙法}で「p 値」を 1 分以内に説明せよ。要求は「\textbf{内在的説明・外在的説明・検出力}」の 3 点を含むこと。\textbf{自己採点}で採れ。
\end{quote}

\textbf{解答の指針:} 内在的説明は「\(H_0\) を真としたとき、観測結果以上に極端なデータが出る確率」。外在的説明の誤りは「\(H_0\) が正しい確率」で、方向が逆。検出力は「\(H_0\) を採用・棄却する基準が有意水準 \(\alpha\)」と対になる量。\textbf{3 点すべて自力で書けたなら合格}。\textbf{外在的説明が書けない}のは、\(H_0\) が棄却される瞬間を一度も理解していないサインです。

\subsection{1 日 1 ファイル・日付名・タグ}

運用は 3 ルールです。

\begin{enumerate}
\item \textbf{1 日 1 ファイル}: ファイル名を日付にする。1 日のバックリンクが「1 枚」に集まり、\textbf{時間順の索引}が自動でできる
\item \textbf{名前には「日付 + 主張」}: \texttt{2026-09-29-p値.md} のように、\textbf{ファイル名に答えを書く}。ファイル名の抽象度を上げない
\item \textbf{タグは 2 個まで}: \texttt{\#memory \#testing}。\textbf{タグが増えるのは失敗のサイン}。\end{enumerate}

\begin{quote}
\textbf{【補足】}\ PARA 系のツールで「観測(observation)」という箱を設ければ、\textbf{道具 1 の文献ノート = Resources の観察記録}に一致します。この 1 箱の分離が、\textbf{「他人の主張」}と\textbf{「自分の変換済みメモ」}を混同しないための具体的な仕組みです。
\end{quote}

\section{道具 3: 作り変える — 自分の形に変える}

\subsection{「変換」とは何か}

道具 3 の中心は、\textbf{入力(他人の言葉)を自分の形に変える}という行為です。\textbf{7 つの技法}を 1 つずつ、具体的な例で示します。

\begin{center}
\begin{tabular}{|l|l|}
\hline
技法 & 意味 \\
\hline
換言 & 同じことを別の言葉で言う \\
\hline
圧縮 & 要約する(3 行に圧縮) \\
\hline
構造化 & 概念どうしの関係を表にする \\
\hline
例示 & 具体例を 1 つ作る \\
\hline
逆算 & 手順を逆にたどる \\
\hline
類比 & 別の分野の既知の言葉で言う \\
\hline
非例 & 例外の例を 1 つ作る \\
\hline
\end{tabular}
\end{center}

\subsection{7 つの技法の例 — 統計学と線形代数から}

\textbf{換言}: 「p 値は帰無仮説のもとでの確率」\(\rightarrow\)「\textbf{観測より希少なデータがもし出ているなら、それはどのくらいの珍しい出来か}」と 1 つの日本語文に訳す。

\textbf{圧縮}: 「相関係数は \(-1 \le r \le 1\) をとり、\(r=0\) は『直線的な関係がない』ことを意味する」——この 2 文を「\textbf{\(r\) の符号が向きを、\textbf{絶対値が直線性の強さ}を表す」の 1 行}に圧縮する。\textbf{1 行にできないなら、まだ理解していない}。

\textbf{構造化}: 「標本平均は母平均の不偏推定量である」という主張を、\textbf{不偏性と頑健性の対比表}にする。

\begin{center}
\begin{tabular}{|l|l|l|}
\hline
推定量 & 不偏性 & 頑健性 \\
\hline
標本平均 & ある(母平均に一致) & 弱い(外れ値に影響される) \\
\hline
標本中央値 & ない(ただし符号は一致) & 強い \\
\hline
\end{tabular}
\end{center}

この表を見て「\textbf{母平均の不偏推定量が唯一の正解か?}」と問うこと自体が、\textbf{構造化}の目的です。

\textbf{例示}: 「\textbf{相関は因果ではない}」に具体例を 1 つ立てる。「アイスクリームの売上と水難事故件数」が典型です(『統計学入門』第 2 節)。

\textbf{逆算}: 「線形写像 \(A\) が可逆とは、\(A\boldsymbol{x} = \boldsymbol{b}\) がすべての \(\boldsymbol{b}\) に対して解をもつこと」——\textbf{手順を逆にたどる}。「\(\boldsymbol{b}\) を自由に選べるなら、\(A\) は押し潰していない」と気づく。\textbf{\(\det A \neq 0\) は「潰さない写像」の条件}になる(『大学数学入門』第 4 節)。

\textbf{類比}: 「作業記憶の容量 \(4 \pm 1\)」に類比をかける——「\textbf{頭で同時に運べる箱が 4 個}」。既知の言葉(箱)で、未知の言葉(容量)を押さえられる。

\textbf{非例}: 「\textbf{中心極限定理には条件がある}」の境目として「\textbf{母集団が重尾なら、標本数を増やしても平均の分布は正規に近づかない}」。

\subsection{Feynman 手法}

\begin{quote}
\textbf{【定義】}\ \textbf{Feynman 手法}とは、\textbf{概念を自分の言葉で説明できないなら、その概念は理解していない}と見なす方法。説明できない箇所が、そのまま未理解の地図になります。
\end{quote}

\begin{quote}
\textbf{【例】}\ 線形代数の「基底」を、自分の言葉で 1 行にすると:「\textbf{互いに独立なベクトルの集合で、それらが張る空間をすべて表せる}」。この 1 行を 2 つの例(\(\mathbb{R}^2\) の標準基底、\(\mathbb{R}^3\) の外積で得た基底)で確認すれば、\textbf{基底とは「独立」+「全域を張る」の二重構造}が見えてきます。
\end{quote}

\begin{quote}
\textbf{【演習】}\ 「\textbf{統計的有意}」を、\textbf{自分の言葉}で 2 行以内で説明せよ。ただし、\textbf{専門語を 1 つも使わない}こと。
\end{quote}

\textbf{解答の指針:} 解答例:「\textbf{偶然なら出るはずだったのに、実際に出た}という運の悪さを、\textbf{5 \% の確率で}が起きたのかどうかを数える」。要点は \textbf{「運が悪かった」の確率を数えている}という一点です(母集団の確率ではなく、標本の確率)。\textbf{専門語を使った説明}は、手順をパスしていません。

\section{道具 3 の実装 — 説明テスト}

\subsection{説明テストの作法}

説明テストは、\textbf{紙 1 枚に新しい概念を書いて、それを初学者に説明させる}という手順です。\textbf{3 人称}に分けると明快です。

\begin{enumerate}
\item \textbf{3 人の想定読者}: ①その分野に入ったばかりの人 ②少し違う分野の人 ③その分野に詳しい人
\item \textbf{3 段階のテスト}: ①一度の説明で通す ②「例は?」で再質問する ③「そうか!」が返るまで繰り返す
\end{enumerate}

\begin{quote}
\textbf{【補足】}\ ③ の詳しい人に通らないなら、\textbf{論理の穴か、暗黙の前提の未説明}です。逆に ① は通っても ③ は通らないなら、\textbf{定義の境界があいまい}です。
\end{quote}

\subsection{説明の錯覚 — illusion of explanatory depth}

説明テストが本当に必要になるのは、\textbf{Rozenblit \& Keil(2002)} の「\textbf{説明の深さの錯覚}」があるからです。\textbf{それらしい説明}には、\textbf{理解とよく似た感触}があります。

\begin{center}
\begin{tabular}{|l|l|}
\hline
項目 & 典型的な水準 \\
\hline
「説明できる」と思う程度 & 高い \\
\hline
実際の説明の深さ & 浅い \\
\hline
\end{tabular}
\end{center}

\begin{quote}
\textbf{【定義】}\ \textbf{説明の深さの錯覚}(illusion of explanatory depth)とは、\textbf{自分が不十分な概念を説明できると思い込んでいる状態}。「\textbf{通じる感じ}」と「\textbf{論理的に正しそう}」を同一視してしまう誤りです。
\end{quote}

\begin{quote}
\textbf{【例】}\ 「相関係数」の説明が「\textbf{相関がある = 関係がある}」で止まっているなら、\textbf{制御変数}(第 3 の要因)を無視した宣言です。「\textbf{共通の要因に両方が動かされている}」ことまで説明できれば、相関と因果の差が言えます。
\end{quote}

\subsection{説明テストの 4 段(実行手順)}

\begin{enumerate}
\item \textbf{定義を自分の言葉で}: 用語を使わずに、上から 2 行
\item \textbf{例を 1 つ}: 具体的な数値で
\item \textbf{非例を 1 つ}: 「\textbf{これは違う}」という境目の例
\item \textbf{他分野との類比}: \textbf{1 行}で、別の既知の現象に翻訳
\end{enumerate}

\begin{quote}
\textbf{【演習】}\ 「\textbf{固有値}」について、説明テストの 4 段(定義・例・非例・類比)を\ \textbf{それぞれ 1 行}で書いてみよ。
\end{quote}

\textbf{解答の指針:} 定義は「\textbf{行列の変換で向きが変わらない特別な方向と、その伸び率}」。例は「\textbf{90 度回転する行列の固有値は複素数で、方向が 90 度回る}」。非例は「\textbf{\(\begin{pmatrix} 1 & 1 \\ 0 & 1 \end{pmatrix}\) は固有値が 2 つとも 1 で、固有ベクトルが無限にある}」こと(重複固有値)。類比は「\textbf{写像の不変な部分だけを取り出す}のと同じ」。\textbf{非例を書けない}なら、4 段の理解はまだ足りません。

\section{道具 4: 引き出す — 隠して思い出す}

\subsection{見返すより、まず隠す}

道具 4 は、\textbf{ノートを隠して自分から引き出す}操作です。\textbf{符号化}が保存を支え、\textbf{検索経路の強度}が使える力を決めます。

\begin{center}
\begin{tabular}{|l|l|}
\hline
方法 & 何をするか \\
\hline
見返す (re-reading) & ノートよく読む(引き出しではない) \\
\hline
検索練習 (retrieval practice) & \textbf{ノートを隠して、回答を先に出してから照合する} \\
\hline
\end{tabular}
\end{center}

\textbf{「見返す」対「隠して思い出す」}——これが本書の最終的な差です。\textbf{『総合問題集』}の「解けなかった問題は 1 週間後に再解く」が、具体的な操作例です。

\subsection{効果量という評価}

\textbf{Dunlosky ほか(2013)} は、学習技術を\textbf{高・中・低}の 3 段階で評価しました。\textbf{検索練習}(practice testing)と\textbf{間隔反復}(distributed practice)が「\textbf{高}」の 2 つで、他の大多数は「低」でした。

\begin{center}
\begin{tabular}{|l|l|}
\hline
技術 & 評価(Dunlosky 2013) \\
\hline
検索練習 & 高 \\
\hline
間隔反復 & 高 \\
\hline
間隔のない再読 & 低 \\
\hline
ハイライト & 低 \\
\hline
要約 & 低 \\
\hline
\end{tabular}
\end{center}

注目すべきは\textbf{「要約」が「低」}である点です。\textbf{道具 3(作り変える)の要約は「低」}、\textbf{道具 4(検索練習)は「高」}。\textbf{「作っただけで済ませる}」のでなく、\textbf{「隠して引き出す}」ので初めて、使える力になります。

\begin{quote}
\textbf{【補足】}\ 『統計学入門』第 12 節でも、\textbf{「書いて動かして間違って直す」}という一連が鍵でした。\textbf{4 道具のサイクル}(収集・構造・変換・出力)は、その実験の記録とも相性がよい。
\end{quote}

\subsection{3 つの引き出し方}

\begin{enumerate}
\item \textbf{質問カード(Anki 系)}: 表面は\textbf{質問}、裏面は\textbf{回答}。\textbf{間隔}を空けて再出題
\item \textbf{自己テスト(紙)}: \textbf{ノートを閉じて}、白紙に要点を\textbf{自分の言葉で}書く。照合は\textbf{後}に
\item \textbf{白紙法(Blank page)}: \textbf{ノートを開かず}、テーマだけから\textbf{自由想起で}書いてから、\textbf{あとにノートを開く}。\textbf{差分がそのまま穴になる}
\end{enumerate}

\begin{quote}
\textbf{【演習】}\ 白紙法で 3 分間、「標準化」のテーマで書き出すとしよ。\textbf{いい運用}と\textbf{悪い運用}を 1 つずつ挙げよ。
\end{quote}

\textbf{解答の指針:} \textbf{いい方}: ノートを閉じた状態で「平均が 0、標準偏差が 1」だけに書いてからノートを開く。\textbf{悪い方}: ノートを開いて\textbf{1 頁目だけ}読んでから書く。これが「\textbf{見返す}」で、\textbf{白紙法}ではない。\textbf{白紙法は「ノートを開く前の状態」がポイント}で、開いた瞬間に終わります。\textbf{テーマだけを用意して 30 秒でも答えさせる}のが現実的です。

\section{道具 4 の実装 — 1 日に 30 分}

\subsection{30 分の日課}

\begin{center}
\begin{tabular}{|l|l|l|}
\hline
時間帯 & 時間 & やること \\
\hline
朝 & 10 分 & \textbf{白紙法}: 昨日のノートを隠して書き出す \\
\hline
夕方 & 20 分 & \textbf{Anki 系}: 昨日追加したカードを 1 回目の抽出 \\
\hline
\end{tabular}
\end{center}

\textbf{この 30 分}は、\textbf{新しいノートを作らない日課}です。\textbf{作る}のは 1 日の他の時間の仕事(道具 1 と 2 に任せます)。\textbf{作らないから 30 分で終わります}。

\subsection{カードの書き方(Anki 系)}

\begin{lstlisting}[language=Text]
表面: p 値は何の確率か(1 行で)
裏面: 帰無仮説 H0 が真のもとで、観測結果以上に極端なデータが
      得られる確率。H0 の正しさの確率ではない(方向が逆)。
      補足: p 値が小さいほど H0 を棄却しやすい
タグ: #testing #統計
間隔: 3 日後に再出題
\end{lstlisting}

\begin{quote}
\textbf{【補足】}\ \textbf{表面が 1 行(質問)}、\textbf{裏面が 2〜4 行(回答 + 落とし穴)}に収まっていること。\textbf{裏面が長くなる}のは、\textbf{表面が広すぎる}サインです。
\end{quote}

\subsection{月例レビュー}

毎日続けても行き詰まります。\textbf{月に 1 度、30 分のレビュー}を行います。

\begin{enumerate}
\item \textbf{1 分〜10 分}: 自分の 1 週間分のファイルを引き出し、\textbf{その週の「不明点」を引き出す}
\item \textbf{10 分〜20 分}: 解決済みの「不明点」に 1 行の回答を追記する
\item \textbf{20 分〜30 分}: \textbf{今月未解決の「不明点」を 3 つだけ}選び、\textbf{今月中に解ける問題}に差し替える
\end{enumerate}

\begin{lstlisting}[language=Text]
2026-09-22 不明点: 母集団が重尾でも中心極限定理が成り立つか
→ 2026-09-29 解決: 成り立たない。標本数を増やしても平均の分布は
  重尾に近づき続ける(非現実的な例: Cauchy 分布の標本平均)
→ 変換したリンク: 2026-09-29-0310([[標本平均]])
\end{lstlisting}

\begin{quote}
\textbf{【演習】}\ \textbf{1 週間だけ 30 分の日課}を実行し、\textbf{1 週間目}と\textbf{2 週間目}の白紙法の記述量を比較せよ。\textbf{差が 0}だった場合の診断を 2 つ挙げよ。
\end{quote}

\textbf{解答の指針:} 診断は 2 つです。①「\textbf{白紙のテーマが曖昧}」で何も出てこない。白紙法には、テーマが 1 つに決まっていることが必要 ②「\textbf{前日にテーマが決まっていない}」ので検索対象が広すぎた。対策は 1 日目だけ、\textbf{「昨日のノートから 1 テーマ選んで白紙法」}までこの 1 段階に絞ることです。\textbf{テーマが狭いほど白紙法は機能}します。\textbf{2 週目}には 5 分でも 3 行でも書けるようになります。

\section{4 道具の統合ワークフロー}

\subsection{1 サイクルの設計}

道具は 1 つずつ独立ではありません。\textbf{1 冊を読む}ときの典型的なサイクルは次の 4 段階です。

\begin{center}
\begin{tabular}{|l|l|l|}
\hline
段階 & 時間 & なにをするか \\
\hline
読む & 2 時間 & 1 日 1 節 \\
\hline
(1) 書き写す & 10 分 & 日付ファイルに 3 つの主張を\textbf{自分の言葉}で残す \\
\hline
(2) 結びつける & 20 分 & 3 つの主張を\textbf{既存のカード}にリンク \\
\hline
(3) 作り変える & 20 分 & \textbf{圧縮}と\textbf{逆算}を 1 つずつ \\
\hline
(4) 引き出す & 翌日 10 分 & \textbf{白紙法}: テーマを 10 行で \\
\hline
\end{tabular}
\end{center}

鍵は、\textbf{「書き写す」と「作り変える」は同じ日にやる}一方、\textbf{「引き出す」は翌日}に回すことです。\textbf{同じ日にやると記憶が温かすぎて、検索練習の効果が落ちます}。\textbf{日をまたぐ}ことが、このサイクルの効きの出どころです。

\begin{quote}
\textbf{【補足】}\ 2 時間の「読む」の中のノート作成は\textbf{合計 50 分}。\textbf{1 日 50 分}で回せます。\textbf{逆に、読む 2 時間を削ってもノート作成は削らない}設計が基本です。
\end{quote}

\subsection{1 冊のノート設計(『統計学入門』の例)}

「\textbf{『統計学入門』を 1 冊処理する}」を具体的な例にします。

\begin{center}
\begin{tabular}{|l|l|l|}
\hline
段階 & 時間 & なにをするか \\
\hline
読む & 2 時間 & 1 日 1 節。『第 1 節 データと統計の考え方』から『第 2 節 1 変量データの要約』まで \\
\hline
(1) 書き写す & 10 分 & 日付ファイルに 3 つの主張を\textbf{自分の言葉}で残す \\
\hline
(2) 結びつける & 20 分 & 3 つの主張を\textbf{他の節}の Zettel カードにリンク \\
\hline
(3) 作り変える & 20 分 & \textbf{圧縮}と\textbf{逆算}を 1 つずつ。4 段の説明テストを 1 つ \\
\hline
(4) 引き出す & 翌日 10 分 & \textbf{白紙法}: 「\textbf{分散と標準偏差}」をテーマ 10 行で \\
\hline
\end{tabular}
\end{center}

これは\textbf{『総合問題集』の設計}とも一致します。\textbf{1 日 1 ファイル}・\textbf{1 主張 1 カード}・\textbf{1 テーマ 1 白紙}——の 3 単位が、そのまま設計の単位です。

\subsection{時間が無い人の縮退版}

すべての時間がないときは、\textbf{1 サイクルで 3 つだけ}できる最小をデザインします。

\begin{enumerate}
\item \textbf{最低 5 分}: 1 日 1 ファイルに\textbf{1 行だけ}書く(\textbf{自分の言葉}になっていること)
\item \textbf{最低 10 分}: その 1 行を、\textbf{昨日のファイル}にリンクする
\item \textbf{最低 10 分}: 昨日の 1 行を\textbf{白紙法}で書き出す
\end{enumerate}

\textbf{道具 1 と道具 4 だけ}なら、\textbf{1 冊 15 分}でサイクルを回せます。この 2 つが押さえられていれば、最低限の合格です。

\begin{quote}
\textbf{【演習】}\ 自分の手元にある本 1 冊(または 1 章)で、\textbf{上記の縮退版を 3 日続けて}実行し、\textbf{3 日目の白紙法で 3 行以上書けるか}を判定せよ。
\end{quote}

\textbf{解答の指針:} 合格の判定は「\textbf{3 行以上、実際の主張(本の一部ではなく自分の言葉)を書けるか}」です。\textbf{3 行目が本文をそのまま写しただけ}なら不合格。\textbf{自分の言葉に変換できていれば、道具 1 は自動判定済}です。\textbf{3 日目には 1 行目がもっとも良い 3 行目}になります。\textbf{本ではなく「あの考え方が少しわかった」という感覚}が 3 日目に出るかどうか——これが唯一のしるしです。

\section{よくある失敗と、その対処}

\subsection{(a) ノートが集まっただけ}

症状: \textbf{ノートは増えたのに、記憶が増えていない}。

原因: \textbf{収集で終わっている}(道具 1 止まり)。

処方: \textbf{1 文で「後の使い道」を書く}。文献ノートに「\textbf{この情報をいつ何に使うか}」を 1 行追加する。\textbf{それを書けないノートは消してよい}。

\subsection{(b) ツールが多すぎる}

症状: \textbf{ノート用のアプリが 10 個ある}。どれを常用すればよいか、分からない。

原因: \textbf{道具 2 の「網」が多系統ある}ため、辿り方が決まらない。

処方: \textbf{1 つのプラットフォームに落ち着く}。\textbf{1 系統のリンクが毎日出る}ようにする。\textbf{ツールの選定基準は「1 年後も使えるか」だけ}。ほかは判断材料にしません。

\subsection{(c) ノートが美しくて中身がない}

症状: \textbf{ノートはきれいなのに、説明できる 1 行も無い}。

原因: \textbf{外在的負荷}に労力を奪われている(第 3 節)。

処方: \textbf{プレーンテキストが原則}。\textbf{装飾してから書く}のではなく、\textbf{テキストで書いてから装飾する}。\textbf{体裁を整えるコストが、理解のコストを上回る}。装飾ゼロでも機能するノートが、\textbf{良いノート}です。

\subsection{(d) 「まとめ」が自己満足}

症状: \textbf{きれいなまとめを 1 度読んで、満足して終わる}。

原因: \textbf{説明の錯覚}。\textbf{「通じる感じ」}だけで満足している。

処方: \textbf{説明テスト}を習慣のなかに加える。\textbf{白紙法で書けなければ、理解はゼロ}として扱う。

\subsection{(e) 消化しないまま放置}

症状: \textbf{「不明点」のノートが 1 年後も放置}。対象が古くなり、あとから見ても意味が通らなくなっている。

原因: 「\textbf{今度}」という言葉は、\textbf{時間}として数えられない。

処方: \textbf{3 ステップで回す}: \textbf{引く(View) → 足す(追加) → 名前を変える(結論に置き換える)}。\textbf{月例レビュー}(第 11 節)でこの 3 ステップを回す。\textbf{それでも解けないものは削除する}。

\begin{quote}
\textbf{【書き手の心構え】}\ ノートの質は文字数ではなく、\textbf{使われかた}で決まる。\textbf{4 つの道具を 1 つずつ回し、最後に白紙に 3 行書ける分だけが残る——それが「将来の自分が引き出せるノート」の全部}だ。残りは、消えても構わない。
\end{quote}

\section{おわりに — 4 行でまとめる}

本書の 4 道具を 1 行ずつに要約します。

\begin{center}
\begin{tabular}{|l|l|}
\hline
道具 & 一行要約 \\
\hline
1 書き写す & 頭のままでは消える。外部に出すと残る。 \\
\hline
2 結びつける & 孤立した知識は検索できない。網を張ると検索できる。 \\
\hline
3 作り変える & 自分の言葉にできないなら、理解していない。 \\
\hline
4 引き出す & 隠して思い出す。見返すのは検索ではない。 \\
\hline
\end{tabular}
\end{center}

そして最後に、この 4 行を\textbf{白紙に書いてください}。\textbf{4 行を書けないなら、まだ本を読んだだけです}。\textbf{「読んだ」で終わらせず、「書いた」にしてください}。

\end{document}
